\chapter{置換}
    \section{置換の符号の定理}
        \subsection{\texorpdfstring{$\sigma'(i) \coloneqq \sigma(n+1-i)$としたときの$\sgn{\sigma'}$}{σ'(i) := σ(n+1-i)としたときのsgn(σ')}}
            \label{置換の引数を逆順にしたときの符号}
            \begin{shadebox}
                $\sigma$を$S=\{1,\dotsc,n\}$上の置換とする。
                $S$上の置換$\sigma'$を$\sigma'(i) = \sigma(n+1-i)$で定めるとき、$\sgn{\sigma'}$は次式となる。
                \begin{align*}
                    \sgn{\sigma'} &= s\sgn{\sigma} \\
                    s &=
                    \begin{cases}
                        1 & (n \text{ is even},\; n/2 \text{ is even}) \\
                        -1 & (n \text{ is even},\; n/2 \text{ is odd}) \\
                        1 & (n \text{ is odd},\; (n-1)/2 \text{ is even}) \\
                        -1 & (n \text{ is odd},\; (n-1)/2 \text{ is odd})
                    \end{cases} \\
                    &=
                    \begin{cases}
                        1 & (n(n-1)/2 \text{ is even}) \\
                        -1 & (n(n-1)/2 \text{ is odd})
                    \end{cases} \\
                    &= (-1)^{n(n-1)/2}
                \end{align*}
            \end{shadebox}
            \begin{proof}
                \quad\par
                Cauchyの2行記法に於いて$\sigma'$の下行は$\sigma$のそれを左右反転したものになっている。
                つまり$\sigma'$を作用させることは$\sigma$を作用させた後に逆順にする置換を施すことに等しい。
                逆順にする置換は、$1,2,\dotsc,n$を左右の端同士で交換する処理を外側から内側に向かって進める方針で考えれば定理の主張が成り立つことがわかる。
            \end{proof}
            Mathematicaを用いた数値例による検証は\verb|theorem_sign_of_permutation.nb|を参照
        \subsection{\texorpdfstring{$\sigma'(i) \coloneqq n+1-\sigma(i)$としたときの$\sgn{\sigma'}$}{σ'(i) := n+1-σ(i)おしたときのsgn(σ')}}
            \label{出力を逆順にしたときの置換の符号}
            \begin{shadebox}
                $\sigma$を$S=\{1,\dotsc,n\}$上の置換とする。
                $S$上の置換$\sigma'$を$\sigma'(i) = n+1-\sigma(i)$で定めるとき、$\sgn{\sigma'}$は次式となる。
                \begin{align*}
                    \sgn{\sigma'} &= s\sgn{\sigma}
                \end{align*}
                $s$は\ref{置換の引数を逆順にしたときの符号}のものと等しい。
            \end{shadebox}
            \begin{proof}
                \quad\par
                $\sigma$を互換の積で表したものを$\sigma=b_k b_{k-1}\dotsm b_1 b_0$とする。
                但し$b_0 \coloneqq \epsilon$(単位置換)とする。
                $\sigma_0 \coloneqq b_0,\;\sigma_l \coloneqq b_l\sigma_{l-1}\;(l=1,2,\dotsc,k)$とする。
                \par
                $\sigma'_0$は逆順にする置換である。
                $1,2,\dotsc,n$を左右の端同士で交換する処理を外側から内側に向かって進める方針で考えれば$\sigma'_0$の符号は定理の主張の$s$に等しいことが解る。
                \par
                $\sigma'_l$と$\sigma'_{l-1}$を比較すると、$b_l$に対応する部分が交換されているので$\sgn{\sigma'_l} = -\sgn{\sigma'_{l-1}}$である。
                これより$\sgn{\sigma'} = \sgn{\sigma'_k} = (-1)^k\sgn{\sigma'_0} = \sgn{\sigma} s$
            \end{proof}
            Mathematicaを用いた数値例による検証は\verb|theorem_sign_of_permutation.nb|を参照